% ============================================================
%  PREAMBULE  -  Business Plan BUT Informatique S5/S6
%  Ne pas ecrire de contenu ici : uniquement reglages et styles
% ============================================================

\usepackage[T1]{fontenc}
\usepackage[utf8]{inputenc}
\usepackage[french]{babel}
\usepackage{lmodern}
\usepackage{microtype}

% ---------- Mise en page ----------
\usepackage[a4paper,top=2.5cm,bottom=2.5cm,left=2.5cm,right=2.5cm]{geometry}
\usepackage{setspace}
\onehalfspacing
\setlength{\parskip}{0.6em}
\setlength{\parindent}{0pt}

% ---------- Couleurs ----------
\usepackage[table]{xcolor}
\definecolor{bpbleu}{HTML}{1F3A5F}     % titres principaux
\definecolor{bpaccent}{HTML}{2E6DA4}   % sous-titres
\definecolor{bpgris}{HTML}{F2F4F7}     % fond d'en-tete de tableau
\definecolor{bpvert}{HTML}{2E7D52}     % forces / opportunites
\definecolor{bprouge}{HTML}{B3261E}    % faiblesses / menaces

% ---------- Titres ----------
\usepackage{titlesec}
\titleformat{\section}{\Large\bfseries\color{bpbleu}}{\thesection.}{0.6em}{}
\titleformat{\subsection}{\large\bfseries\color{bpaccent}}{\thesubsection}{0.6em}{}
\titleformat{\subsubsection}{\normalsize\bfseries}{\thesubsubsection}{0.6em}{}
\titlespacing*{\section}{0pt}{1.6em}{0.8em}

% ---------- En-tetes et pieds de page ----------
\usepackage{fancyhdr}
\pagestyle{fancy}
\fancyhf{}
\fancyhead[L]{\small\color{bpbleu}\nomentreprise}
\fancyhead[R]{\small\color{bpbleu}Business Plan \textendash{} \anneeuniv}
\fancyfoot[C]{\small\thepage}
\renewcommand{\headrulewidth}{0.4pt}

% ---------- Tableaux ----------
\usepackage{booktabs}
\usepackage{tabularx}
\usepackage{longtable}
\usepackage{array}
\usepackage{multirow}
\usepackage{ragged2e}
\renewcommand{\arraystretch}{1.35}

% Colonne X justifiee et colonne X centree
\newcolumntype{L}{>{\RaggedRight\arraybackslash}X}
\newcolumntype{C}{>{\Centering\arraybackslash}X}
\newcolumntype{P}[1]{>{\RaggedRight\arraybackslash}p{#1}}

% ---------- Images, listes, liens ----------
\usepackage{graphicx}
\graphicspath{{img/}}
\usepackage{enumitem}
\setlist[itemize]{leftmargin=1.4em,itemsep=0.2em,topsep=0.3em}
\usepackage{caption}
\captionsetup{font=small,labelfont={bf,color=bpbleu}}
\usepackage[hidelinks]{hyperref}
\usepackage{url}
\urlstyle{same}

% ============================================================
%  COMMANDES DU PROJET  -  a modifier une seule fois
% ============================================================
\newcommand{\nomentreprise}{NOM\_ENTREPRISE}
\newcommand{\anneeuniv}{2026\,--\,2027}
\newcommand{\etudiantun}{Antoine Barthélémy}
\newcommand{\etudiantdeux}{Simisola Agbe}
\newcommand{\cycle}{S5/S6 BUT Informatique \textendash{} IUT de Reims-Châlons-Charleville}
\newcommand{\prof}{Gilles Escuyer}
\newcommand{\profmail}{Gilles.Escuyer@univ-reims.fr}

% ============================================================
%  ENVIRONNEMENTS REUTILISABLES
% ============================================================

% --- En-tete de tableau grise et en gras ---
\newcommand{\thead}[1]{\cellcolor{bpgris}\textbf{#1}}

% --- Encadre "hypothese a valider" ---
\usepackage[most]{tcolorbox}
\newtcolorbox{hypothese}{
  colback=bpgris, colframe=bpaccent, boxrule=0.6pt,
  left=8pt, right=8pt, top=6pt, bottom=6pt,
  title=\textbf{Hypothèse à valider}, fonttitle=\small\color{white},
  colbacktitle=bpaccent
}

% --- Tableau avantages / inconvenients (etape I) ---
% Usage : \begin{avantinconv}{Titre de l'idee}
%           Avantage & Inconvenient \\
%         \end{avantinconv}
\newenvironment{avantinconv}[1]
{\par\vspace{0.4em}\noindent\textbf{#1}\par\vspace{0.3em}
 \tabularx{\textwidth}{LL}
 \toprule
 \thead{Avantages} & \thead{Inconvénients} \\
 \midrule}
{\bottomrule\endtabularx\par\vspace{0.8em}}
